% ===================================================================
%  Подготовка к контрольной работе по семинарам 1–4.
%  Содержательная часть — единый источник для обоих файлов.
%
%  Формат задач — по заданиям 2024/25 («Семинар 5.docx», «Семинар 5
%  (приложение).xlsx»): теоретические вопросы (смещение и состоятельность,
%  R^2_adj, AIC/BIC, мультиколлинеарность), VIF по вспомогательным R^2,
%  парная регрессия по таблице «вручную» с выводом формул МНК.
%  Остальные задачи — из «Задачника на сайт ВШЭ», части 2–4: темы
%  семинаров 1–4, задачи, которые на семинарах не разбирались.
%
%  ОБОЗНАЧЕНИЯ. Как во всех семинарах: RSS — объяснённая сумма квадратов,
%  ESS — остаточная. В задачнике ВШЭ наоборот (RSS — остаточная), поэтому
%  условия переписаны в наших обозначениях, и в каждом условии словами
%  сказано, какая сумма имеется в виду.
%
%  Раздел «Основные формулы» и всё внутри \begin{solution}...\end{solution}
%  попадает только в kr_1_solutions.pdf.
%  Все числа пересчитаны независимо; задача 1 совпадает с выводом Пакета
%  анализа в «Семинар 5 (приложение).xlsx».
% ===================================================================

% «Вопрос N» — оформление как у \problem
\newcommand{\question}[1]{%
  \bigskip\par\noindent
  {\bfseries\color{accent}\large Вопрос #1}\par\nopagebreak\smallskip}

\seminartitle{Подготовка к контрольной}{Семинары 1--4: парная и множественная регрессия,\\
фиктивные переменные, мультиколлинеарность}

% ===================================================================
\ifsolutions
\section{Основные формулы}

\subsection{Парная регрессия}\label{sec:pair}

Модель $Y_t = \alpha + \beta X_t + \varepsilon_t$, $t = 1, \dots, n$. Суммы в
отклонениях ($x_t = X_t - \overline{X}$, $y_t = Y_t - \overline{Y}$):
\begin{gather}
  \sum x_t^2 = \sum X_t^2 - \frac{(\sum X_t)^2}{n},
  \qquad
  \sum y_t^2 = \sum Y_t^2 - \frac{(\sum Y_t)^2}{n},
  \label{eq:dev} \\[2pt]
  \sum x_ty_t = \sum X_tY_t - \frac{\sum X_t\sum Y_t}{n}.
  \notag
\end{gather}
МНК-оценки, суммы квадратов и коэффициент детерминации:
\begin{gather}
  \widehat{\beta} = \frac{\sum x_ty_t}{\sum x_t^2},
  \qquad
  \widehat{\alpha} = \overline{Y} - \widehat{\beta}\,\overline{X},
  \label{eq:ols} \\[2pt]
  \RSS = \widehat{\beta}^2\sum x_t^2,
  \qquad
  \ESS = \sum y_t^2 - \RSS,
  \qquad
  R^2 = \frac{\RSS}{\TSS} = 1 - \frac{\ESS}{\TSS}.
  \notag
\end{gather}
Здесь $\TSS = \sum y_t^2$, $\RSS$ — объяснённая сумма квадратов, $\ESS =
\sum e_t^2$ — остаточная. В парной регрессии $R^2 = r_{XY}^2$.

\subsection{Оценка дисперсии, $t$-тест, доверительный интервал}\label{sec:t}

\begin{gather}
  s^2 = \frac{\ESS}{n-k},
  \qquad
  s_{\widehat{\beta}}^2 = \frac{s^2}{\sum x_t^2}\ \ (\text{парная регрессия}),
  \label{eq:t} \\[2pt]
  t = \frac{\widehat{\beta}_j - \beta_j^0}{s_{\widehat{\beta}_j}} \sim t(n-k),
  \qquad
  \widehat{\beta}_j \pm t_{\alpha/2}(n-k)\,s_{\widehat{\beta}_j},
  \notag
\end{gather}
$k$ — число параметров вместе со свободным членом ($k = 2$ в парной
регрессии). Гипотеза $H_0\colon \beta_j = \beta_j^0$ отвергается, если
$|t| > t_{\alpha/2}(n-k)$.

\subsection{$F$-тесты и скорректированный $R^2$}\label{sec:f}

Значимость регрессии в целом ($H_0$: все коэффициенты, кроме свободного
члена, равны нулю):
\begin{equation}
  F = \frac{\RSS/(k-1)}{\ESS/(n-k)}
    = \frac{R^2}{1-R^2}\cdot\frac{n-k}{k-1} \sim F(k-1,\ n-k).
  \label{eq:f}
\end{equation}
Тест на $q$ линейных ограничений (R — модель с ограничениями, UR — без них,
$k$ — число параметров модели UR):
\begin{equation}
  F = \frac{\bigl(\ESS_{\mathrm{R}} - \ESS_{\mathrm{UR}}\bigr)/q}
           {\ESS_{\mathrm{UR}}/(n-k)} \sim F(q,\ n-k).
  \label{eq:restr}
\end{equation}
Тест Чоу для двух подвыборок: $\ESS_{\mathrm{UR}} = \ESS_1 + \ESS_2$, $q = k$,
знаменатель $\ESS_{\mathrm{UR}}/(n - 2k)$.

Скорректированный коэффициент детерминации и информационные критерии:
\begin{gather}
  R^2_{\text{adj}}
  = 1 - \frac{\ESS/(n-k)}{\TSS/(n-1)}
  = 1 - \frac{n-1}{n-k}\bigl(1 - R^2\bigr),
  \label{eq:adj} \\[2pt]
  \mathrm{AIC} = \frac{2k}{n} + \ln\frac{\ESS}{n},
  \qquad
  \mathrm{BIC} = \frac{k\ln n}{n} + \ln\frac{\ESS}{n}.
  \notag
\end{gather}

\subsection{Мультиколлинеарность}\label{sec:vif}

$R_j^2$ — коэффициент детерминации регрессии $j$-го регрессора на остальные
(с константой):
\begin{equation}
  \mathrm{VIF}_j = \frac{1}{1 - R_j^2},
  \qquad
  \D(\widehat{\beta}_j) = \mathrm{VIF}_j\cdot\frac{\sigma^2}{\sum_t(X_{tj} - \overline{X}_j)^2}.
  \label{eq:vif}
\end{equation}
Пороги: $\mathrm{VIF}_j > 5$ — заметная мультиколлинеарность,
$\mathrm{VIF}_j > 10$ ($R_j^2 > 0{,}9$) — серьёзная. При двух регрессорах
$R_j^2 = r^2$, и оба VIF одинаковы.
\fi

% ===================================================================
\section{Теоретические вопросы}

% -------------------------------------------------------------------
\question{1}

Что такое смещение оценки $\widehat{\beta}$ параметра $\beta$? Какая оценка
называется несмещённой и какая — состоятельной? Приведите пример оценки,
которая смещена, но состоятельна.

\begin{solution}
\emph{Смещение} — разность между математическим ожиданием оценки и истинным
значением параметра:
\[
  \operatorname{bias}(\widehat{\beta}) = \E\widehat{\beta} - \beta .
\]
Оценка \emph{несмещённая}, если $\E\widehat{\beta} = \beta$ при любом объёме
выборки: в среднем по многим выборкам она попадает в истинное значение.

Оценка $\widehat{\beta}_n$ \emph{состоятельна}, если при росте объёма выборки
она сходится по вероятности к истинному значению:
\[
  \lim_{n\to\infty}\Pr\bigl\{|\widehat{\beta}_n - \beta| > \varepsilon\bigr\} = 0
  \quad\text{для любого } \varepsilon > 0 .
\]
По мере роста выборки оценка всё ближе к истине. Достаточное условие:
$\E\widehat{\beta}_n \to \beta$ и $\D(\widehat{\beta}_n) \to 0$.

\emph{Пример.} Оценка дисперсии ошибок $\widehat{\sigma}^2 = \ESS/n$ смещена:
$\E\widehat{\sigma}^2 = \frac{n-k}{n}\sigma^2 < \sigma^2$. Но смещение
$-\frac{k}{n}\sigma^2$ исчезает при $n \to \infty$, и оценка состоятельна.
Несмещённая оценка — $s^2 = \ESS/(n-k)$.
\end{solution}

% -------------------------------------------------------------------
\question{2}

Перечислите предпосылки классической линейной регрессионной модели.
Сформулируйте теорему Гаусса--Маркова. Что означает, что МНК-оценка —
BLUE?

\begin{solution}
Модель $\mathbf{y} = \mathbf{X}\boldsymbol{\beta} + \boldsymbol{\varepsilon}$.
Предпосылки:
\begin{enumerate}[label=\arabic*), leftmargin=2.2em, itemsep=2pt]
  \item модель линейна по параметрам и верно специфицирована;
  \item $\mathbf{X}$ детерминирована, столбцы линейно независимы
        ($\operatorname{rank}\mathbf{X} = k < n$) — нет строгой
        мультиколлинеарности;
  \item $\E\varepsilon_t = 0$;
  \item $\D(\varepsilon_t) = \sigma^2$ для всех $t$ — гомоскедастичность;
  \item $\Cov(\varepsilon_t, \varepsilon_s) = 0$ при $t \ne s$ — нет
        автокорреляции (3–5 в матричном виде: $\E\boldsymbol{\varepsilon} =
        \mathbf{0}$, $\D(\boldsymbol{\varepsilon}) = \sigma^2\mathbf{I}$);
  \item (для тестов) $\varepsilon_t \sim N(0, \sigma^2)$.
\end{enumerate}

\textbf{Теорема Гаусса--Маркова.} При выполнении предпосылок 1–5 МНК-оценка
$\widehat{\boldsymbol{\beta}} = (\mathbf{X}^{\top}\mathbf{X})^{-1}\mathbf{X}^{\top}\mathbf{y}$
имеет наименьшую дисперсию в классе линейных (по $\mathbf{y}$) несмещённых
оценок.

\textbf{BLUE} — Best Linear Unbiased Estimator: \emph{линейная} (линейная
функция от $\mathbf{y}$), \emph{несмещённая} ($\E\widehat{\boldsymbol{\beta}} =
\boldsymbol{\beta}$) и \emph{наилучшая}, то есть эффективная — с наименьшей
дисперсией среди всех линейных несмещённых оценок. Нормальность ошибок для
теоремы не нужна.
\end{solution}

% -------------------------------------------------------------------
\question{3}

Напишите формулу скорректированного коэффициента детерминации и объясните
смысловое отличие от обычного $R^2$. Сформулируйте информационные критерии
Акаике (AIC) и Шварца (BIC): формулы и для чего они используются.

\begin{solution}
По \eqref{eq:adj}
\[
  R^2_{\text{adj}} = 1 - \frac{\ESS/(n-k)}{\TSS/(n-1)}
                   = 1 - \frac{n-1}{n-k}\bigl(1 - R^2\bigr).
\]
Обычный $R^2 = 1 - \ESS/\TSS$ не убывает при добавлении \emph{любого}
регрессора, даже бессмысленного: $\ESS$ может только уменьшиться. Поэтому
сравнивать по $R^2$ модели с разным числом регрессоров нельзя. В
$R^2_{\text{adj}}$ суммы квадратов делятся на свои степени свободы, и
добавление регрессора увеличивает его, только если выигрыш в $\ESS$
перевешивает потерю степени свободы. Всегда $R^2_{\text{adj}} \le R^2$,
$R^2_{\text{adj}}$ может быть отрицательным.

\[
  \mathrm{AIC} = \frac{2k}{n} + \ln\frac{\ESS}{n},
  \qquad
  \mathrm{BIC} = \frac{k\ln n}{n} + \ln\frac{\ESS}{n},
\]
$n$ — число наблюдений, $k$ — число параметров вместе с константой. Критерии
нужны для выбора между моделями (в том числе не вложенными друг в друга) по
одним и тем же данным: второе слагаемое — качество подгонки, первое — штраф
за число параметров. Лучше модель с \emph{меньшим} значением критерия. Штраф
BIC строже ($\ln n > 2$ при $n \ge 8$), поэтому BIC выбирает более экономные
модели.
\end{solution}

% -------------------------------------------------------------------
\question{4}

По одной и той же выборке оценены регрессии $Y$ на $X$ и $X$ на $Y$.
Совпадут ли линии регрессии?

\begin{solution}
Оценим обе регрессии по МНК:
\[
  \text{(1) } Y \text{ на } X:\quad \widehat{Y} = \widehat{a} + \widehat{b}_{Y|X}X,
  \qquad
  \text{(2) } X \text{ на } Y:\quad \widehat{X} = \widehat{c} + \widehat{b}_{X|Y}Y,
\]
\[
  \widehat{b}_{Y|X} = \frac{\sum x_ty_t}{\sum x_t^2},
  \quad
  \widehat{a} = \overline{Y} - \widehat{b}_{Y|X}\overline{X};
  \qquad
  \widehat{b}_{X|Y} = \frac{\sum x_ty_t}{\sum y_t^2},
  \quad
  \widehat{c} = \overline{X} - \widehat{b}_{X|Y}\overline{Y}.
\]
Во второй регрессии $X$ и $Y$ просто поменялись ролями, поэтому её формулы
получаются из формул первой заменой $X \leftrightarrow Y$.

\medskip
\textbf{Шаг 1. Обе прямые проходят через точку средних.} Это следует из
формулы свободного члена — первого нормального уравнения МНК. Подставим
$X = \overline{X}$ в первую прямую:
\[
  \widehat{a} + \widehat{b}_{Y|X}\overline{X}
  = \bigl(\overline{Y} - \widehat{b}_{Y|X}\overline{X}\bigr) + \widehat{b}_{Y|X}\overline{X}
  = \overline{Y}.
\]
Так же для второй: при $Y = \overline{Y}$ получаем
$\widehat{c} + \widehat{b}_{X|Y}\overline{Y} = \overline{X}$. Значит, точка
$(\overline{X}, \overline{Y})$ лежит на обеих прямых.

\medskip
\textbf{Шаг 2. Записываем обе прямые в одних осях.} Сравнивать прямые можно,
только если они нарисованы в одной системе координат: $X$ по горизонтали, $Y$
по вертикали. Первая прямая уже так записана: её наклон (насколько растёт $Y$
при росте $X$ на единицу) равен $\widehat{b}_{Y|X}$. Вторая записана
«наоборот» — $X$ через $Y$. Выразим из неё $Y$:
\[
  X = \widehat{c} + \widehat{b}_{X|Y}Y
  \quad\Longrightarrow\quad
  Y = -\frac{\widehat{c}}{\widehat{b}_{X|Y}} + \frac{1}{\widehat{b}_{X|Y}}\,X .
\]
Отсюда и обратная величина: если по второй прямой $X$ растёт на
$\widehat{b}_{X|Y}$ при росте $Y$ на единицу, то $Y$ растёт на
$1/\widehat{b}_{X|Y}$ при росте $X$ на единицу. В осях $(X, Y)$ наклон второй
прямой равен $1/\widehat{b}_{X|Y}$, а не $\widehat{b}_{X|Y}$.

\medskip
\textbf{Шаг 3. Почему достаточно сравнить наклоны.} Прямая на плоскости
однозначно задаётся одной своей точкой и наклоном. Общая точка у наших прямых
уже есть — $(\overline{X}, \overline{Y})$ из шага 1. Поэтому прямые совпадают
тогда и только тогда, когда совпадают наклоны, а свободные члены отдельно
проверять не нужно: они определяются наклоном и общей точкой и при равных
наклонах совпадут автоматически. Проверим явно. Свободный член первой прямой
$\widehat{a} = \overline{Y} - \widehat{b}_{Y|X}\overline{X}$; второй (из
шага 2)
\[
  -\frac{\widehat{c}}{\widehat{b}_{X|Y}}
  = -\frac{\overline{X} - \widehat{b}_{X|Y}\overline{Y}}{\widehat{b}_{X|Y}}
  = \overline{Y} - \frac{1}{\widehat{b}_{X|Y}}\,\overline{X}.
\]
Оба имеют вид «$\overline{Y}$ минус наклон, умноженный на $\overline{X}$», и
при $\widehat{b}_{Y|X} = 1/\widehat{b}_{X|Y}$ равны.

\medskip
\textbf{Шаг 4. Условие совпадения.}
\[
  \widehat{b}_{Y|X} = \frac{1}{\widehat{b}_{X|Y}}
  \iff
  \widehat{b}_{Y|X}\,\widehat{b}_{X|Y} = 1
  \iff
  \frac{\sum x_ty_t}{\sum x_t^2}\cdot\frac{\sum x_ty_t}{\sum y_t^2}
  = \frac{(\sum x_ty_t)^2}{\sum x_t^2\sum y_t^2} = r^2 = 1 .
\]
Прямые совпадают только при $|r| = 1$, то есть когда все точки лежат на одной
прямой. При $|r| < 1$ они пересекаются в точке средних под углом: так как
$|\widehat{b}_{Y|X}| = r^2/|\widehat{b}_{X|Y}| < 1/|\widehat{b}_{X|Y}|$,
вторая прямая всегда круче первой. В крайнем случае $r = 0$ первая прямая
горизонтальна ($Y = \overline{Y}$), а вторая вертикальна ($X = \overline{X}$).

\begin{center}
\begin{tikzpicture}[>=stealth, line join=round, x=0.8cm, y=0.8cm]
  % 10 точек; средние (4,08; 2,91); b_{Y|X} = 0,306; 1/b_{X|Y} = 0,990; r^2 = 0,31
  \draw[->] (0,0) -- (7.8,0) node[anchor=north] {$X$};
  \draw[->] (0,0) -- (0,5.6) node[anchor=east]  {$Y$};

  \draw[dotted, black!55] (4.08,0) -- (4.08,2.91) -- (0,2.91);
  \node[anchor=north] at (4.08,0) {$\overline{X}$};
  \node[anchor=east]  at (0,2.91) {$\overline{Y}$};

  \draw[very thick, accent] (0.8,1.905) -- (7.2,3.866)
        node[anchor=west, font=\small] {$Y$ на $X$: наклон $\widehat{b}_{Y|X} = 0{,}31$};
  \draw[very thick, solbar] (1.95,0.8) -- (6.19,5.0)
        node[anchor=south, font=\small] {$X$ на $Y$: наклон $1/\widehat{b}_{X|Y} = 0{,}99$};

  \foreach \p in {(1.5,2.6),(2.2,1.4),(2.8,3.1),(3.3,1.9),(3.8,3.9),
                  (4.3,2.3),(4.9,3.7),(5.4,2.6),(6.0,4.3),(6.6,3.3)}
    \fill[black!60] \p circle (1.8pt);
  \fill (4.08,2.91) circle (2.6pt);
\end{tikzpicture}

\smallskip
{\small\itshape Пример по 10 точкам. Обе прямые проходят через точку средних,
но наклоны в осях $(X, Y)$ разные: $0{,}31$ и $0{,}99$. Прямые не совпадают,
потому что $\widehat{b}_{Y|X}\,\widehat{b}_{X|Y} = 0{,}306\cdot 1{,}010 = r^2 = 0{,}31 \ne 1$.\par}
\end{center}

\textbf{Почему так.} Регрессия $Y$ на $X$ минимизирует сумму квадратов
отклонений точек от прямой \emph{по вертикали} (ошибки в $Y$), а регрессия $X$
на $Y$ — \emph{по горизонтали} (ошибки в $X$). Это разные задачи, и при
неидеальной связи у них разные решения.
\end{solution}

% ===================================================================
\ifsolutions\clearpage\fi
\section{Задачи}

% -------------------------------------------------------------------
\problem{1 (задание 2024/25)}

В таблице приведены ежегодные значения денежной массы и национального дохода
некоторой гипотетической страны (в миллиардах денежных единиц).

\begin{center}
\begin{tabular}{ccc @{\qquad} ccc}
\toprule
Год & Денежная масса & Нац. доход & Год & Денежная масса & Нац. доход \\
\midrule
1981 & 2,0 & 5,0 & 1986 & 4,0 & 7,7 \\
1982 & 2,5 & 5,5 & 1987 & 4,2 & 8,4 \\
1983 & 3,2 & 6,0 & 1988 & 4,6 & 9,0 \\
1984 & 3,6 & 7,0 & 1989 & 4,8 & 9,7 \\
1985 & 3,3 & 7,2 & 1990 & 5,0 & 10,0 \\
\bottomrule
\end{tabular}
\end{center}

\begin{enumerate}[label=\asbuk*), leftmargin=2.2em, itemsep=3pt]
  \item Оцените регрессию национального дохода ($Y$) на денежную массу ($X$)
        и константу. Обоснуйте вывод основных формул.
  \item Рассчитайте $R^2$ модели.
  \item Проверьте значимость коэффициента при денежной массе на уровне 5\,\%.
\end{enumerate}

\begin{solution}
\textbf{а) Вывод формул МНК.} МНК выбирает $\alpha$ и $\beta$, минимизирующие
сумму квадратов отклонений
\[
  S(\alpha, \beta) = \sum_{t=1}^{n}\bigl(Y_t - \alpha - \beta X_t\bigr)^2 .
\]
Приравниваем нулю частные производные:
\[
  \frac{\partial S}{\partial\alpha} = -2\sum\bigl(Y_t - \alpha - \beta X_t\bigr) = 0,
  \qquad
  \frac{\partial S}{\partial\beta} = -2\sum X_t\bigl(Y_t - \alpha - \beta X_t\bigr) = 0,
\]
и получаем систему нормальных уравнений
\[
  \begin{cases}
    n\alpha + \beta\sum X_t = \sum Y_t, \\[2pt]
    \alpha\sum X_t + \beta\sum X_t^2 = \sum X_tY_t .
  \end{cases}
\]
Из первого уравнения $\alpha = \overline{Y} - \beta\overline{X}$. Подставляя
во второе, получаем
\[
  \widehat{\beta} = \frac{\sum X_tY_t - n\overline{X}\,\overline{Y}}{\sum X_t^2 - n\overline{X}^2}
                  = \frac{\sum x_ty_t}{\sum x_t^2},
  \qquad
  \widehat{\alpha} = \overline{Y} - \widehat{\beta}\,\overline{X}.
\]
Вторые производные $\partial^2S/\partial\alpha^2 = 2n$ и определитель матрицы
вторых производных $4\bigl(n\sum X_t^2 - (\sum X_t)^2\bigr) = 4n\sum x_t^2 > 0$,
так что это минимум.

\textbf{Расчёт.} $n = 10$. Суммы по таблице:
\[
  \sum X_t = 37{,}2,\quad \sum Y_t = 75{,}5,\quad \sum X_t^2 = 147{,}18,
\]
\[
  \sum Y_t^2 = 597{,}03,\quad \sum X_tY_t = 295{,}95 .
\]
Средние $\overline{X} = 3{,}72$, $\overline{Y} = 7{,}55$. По \eqref{eq:dev}
\[
\begin{aligned}
  \sum x_t^2  &= 147{,}18 - \frac{37{,}2^2}{10} = 147{,}18 - 138{,}384 = 8{,}796, \\
  \sum x_ty_t &= 295{,}95 - \frac{37{,}2\cdot 75{,}5}{10} = 295{,}95 - 280{,}86 = 15{,}09, \\
  \sum y_t^2  &= 597{,}03 - \frac{75{,}5^2}{10} = 597{,}03 - 570{,}025 = 27{,}005 .
\end{aligned}
\]
\[
  \widehat{\beta} = \frac{15{,}09}{8{,}796} = 1{,}716,
  \qquad
  \widehat{\alpha} = 7{,}55 - 1{,}716\cdot 3{,}72 = 1{,}168,
  \qquad
  \widehat{Y} = 1{,}168 + 1{,}716\,X .
\]
Рост денежной массы на 1 млрд сопровождается в среднем ростом
национального дохода на 1,72 млрд.

\medskip
\textbf{б) Коэффициент детерминации.} По \eqref{eq:ols}
\[
  \RSS = \widehat{\beta}^2\sum x_t^2 = 1{,}716^2\cdot 8{,}796 = 25{,}888,
  \qquad
  R^2 = \frac{\RSS}{\TSS} = \frac{25{,}888}{27{,}005} = 0{,}959 .
\]
Денежная масса объясняет 95,9\,\% разброса национального дохода.

\medskip
\textbf{в) Значимость наклона.} $\ESS = 27{,}005 - 25{,}888 = 1{,}117$, и по
\eqref{eq:t}
\[
  s^2 = \frac{\ESS}{n-2} = \frac{1{,}117}{8} = 0{,}1397,
  \qquad
  s_{\widehat{\beta}} = \sqrt{\frac{0{,}1397}{8{,}796}} = 0{,}126,
\]
\[
  t = \frac{1{,}716}{0{,}126} = 13{,}6 > 2{,}306 = t_{0{,}025}(8).
\]
Коэффициент значим.
\end{solution}

% -------------------------------------------------------------------
\problem{2 (задачник ВШЭ, №2.3.2)}

По ежегодным данным за 1975--1988 гг. (14 наблюдений) оценивалась
зависимость цены на бензин ($\mathit{Petrol}$, центы за галлон) от цены на
сырую нефть ($\mathit{Oil}$, доллары за баррель):
\[
  \widehat{\mathit{Petrol}}_i = \underset{(4{,}6)}{41{,}9} + \underset{(0{,}2)}{3{,}0}\,\mathit{Oil}_i ,
\]
в скобках — стандартные ошибки. Сумма квадратов остатков $\ESS = 631{,}1$,
$\TSS = 12622$. Предпосылки классической линейной нормальной модели
выполнены.
\begin{enumerate}[label=\asbuk*), leftmargin=2.2em, itemsep=3pt]
  \item Рассчитайте коэффициент детерминации $R^2$.
  \item Проверьте гипотезу о том, что рост цены на нефть на 1 долл. за баррель
        соответствует росту цены на бензин на 2 цента за галлон, на уровне
        значимости 10\,\%.
  \item В 1988 году цена на нефть составила 12,57 доллара за баррель. Какова
        ожидаемая цена на бензин? Какой должна быть цена на нефть, чтобы
        ожидаемая цена на бензин составила 100 центов за галлон?
\end{enumerate}

\begin{solution}
\textbf{а)} $R^2 = 1 - \dfrac{\ESS}{\TSS} = 1 - \dfrac{631{,}1}{12622} = 0{,}95$.

\medskip
\textbf{б)} $H_0\colon \beta = 2$ против $H_1\colon \beta \ne 2$. По
\eqref{eq:t} с $n - k = 12$:
\[
  t = \frac{\widehat{\beta} - 2}{s_{\widehat{\beta}}} = \frac{3{,}0 - 2}{0{,}2} = 5
  > 1{,}782 = t_{0{,}05}(12).
\]
Гипотеза отвергается: бензин дорожает сильнее, чем на 2 цента за доллар.

\medskip
\textbf{в)} Прогноз: $41{,}9 + 3{,}0\cdot 12{,}57 = 79{,}61$ цента за галлон.
Цена нефти для ожидаемой цены бензина 100 центов:
$41{,}9 + 3{,}0\cdot\mathit{Oil} = 100$, откуда
$\mathit{Oil} = (100 - 41{,}9)/3{,}0 = 19{,}37$ доллара за баррель.
\end{solution}

% -------------------------------------------------------------------
\problem{3 (задачник ВШЭ, №2.3.5)}

Рассчитанный по 20 наблюдениям за признаками $X$ и $Y$ выборочный
коэффициент корреляции составил 0,5. Есть ли основания считать, что между
признаками существует регрессионная зависимость? Сформулируйте основную и
альтернативную гипотезы в терминах коэффициентов регрессии; уровень
значимости 5\,\%.

\begin{solution}
Модель $Y_t = \alpha + \beta X_t + \varepsilon_t$; зависимость есть, если
$\beta \ne 0$: $H_0\colon \beta = 0$, $H_1\colon \beta \ne 0$.

$t$-статистику наклона можно выразить через $r$. Так как
$\widehat{\beta} = \sum x_ty_t/\sum x_t^2$, $\ESS = (1 - r^2)\sum y_t^2$ и
$r^2 = (\sum x_ty_t)^2/(\sum x_t^2\sum y_t^2)$,
\[
  t = \frac{\widehat{\beta}}{s_{\widehat{\beta}}}
    = \frac{\widehat{\beta}\sqrt{\sum x_t^2}}{\sqrt{\ESS/(n-2)}}
    = \frac{r\sqrt{n-2}}{\sqrt{1 - r^2}} .
\]
Здесь
\[
  t = \frac{0{,}5\sqrt{18}}{\sqrt{1 - 0{,}25}} = 2{,}449 > 2{,}101 = t_{0{,}025}(18),
\]
гипотеза $H_0$ отвергается: регрессионная зависимость есть.
\end{solution}

% -------------------------------------------------------------------
\problem{4 (задачник ВШЭ, №3.1.3)}

При оценивании по 40 наблюдениям зависимости
$Y_i = \beta_1 + \beta_2X_{2i} + \beta_3X_{3i} + \beta_4X_{4i} + \varepsilon_i$
получены $R^2 = 0{,}6$ и $\TSS = 200$. Рассчитайте несмещённую оценку
дисперсии случайной составляющей и скорректированный коэффициент
детерминации.

\begin{solution}
$n = 40$, $k = 4$, $n - k = 36$. Остаточная сумма
$\ESS = (1 - R^2)\TSS = 0{,}4\cdot 200 = 80$, и по \eqref{eq:t}, \eqref{eq:adj}
\[
  s^2 = \frac{\ESS}{n-k} = \frac{80}{36} = 2{,}22,
  \qquad
  R^2_{\text{adj}} = 1 - \frac{n-1}{n-k}\bigl(1 - R^2\bigr) = 1 - \frac{39}{36}\cdot 0{,}4 = 0{,}567 .
\]
\end{solution}

% -------------------------------------------------------------------
\problem{5 (задачник ВШЭ, №3.2.4)}

Исследователь оценил регрессию цены подержанного автомобиля:
\[
  \mathit{Цена}_i = \beta_1 + \beta_2\,\mathit{Пробег}_i + \beta_3\,\mathit{Поломки}_i
  + \beta_4\,\mathit{Иномарка}_i + \beta_5\,\mathit{Литраж}_i + \beta_6\,\mathit{Длина}_i
  + \varepsilon_i ,
\]
где $\mathit{Иномарка}_i = 1$ для иномарок и 0 для отечественных машин.
Результаты (Model SS — объяснённая сумма квадратов $\RSS$, Residual SS —
остаточная $\ESS$):

\begin{center}
\small
\begin{tabular}{lrrr @{\qquad} lr}
\toprule
       & SS & df & MS & \multicolumn{2}{l}{} \\
\midrule
Model    & 176\,083\,058 &  5 & 35\,216\,611,6 & Number of obs & 50 \\
Residual & 318\,894\,258 & 44 &  7\,247\,596,78 & $F(5, 44)$ & 4,86 \\
Total    & 494\,977\,316 & 49 & 10\,101\,577,9 & Prob $> F$ & 0,0013 \\
         &               &    &                & $R^2$ & 0,3557 \\
         &               &    &                & Adj $R^2$ & \textbf{XXXX} \\
\bottomrule
\end{tabular}

\medskip
\begin{tabular}{lrrrrrr}
\toprule
 & Coef. & Std. Err. & $t$ & $P>|t|$ & \multicolumn{2}{c}{95\,\% Conf. Interval} \\
\midrule
пробег   & $-321{,}62$ & 172,21 & $-1{,}87$ & 0,068 & $-668{,}68$ & 25,44 \\
поломки  & $-466{,}65$ & 603,61 & $-0{,}77$ & 0,444 & $-1683{,}15$ & 749,84 \\
иномарка & 5102,07 & 1749,39 & 2,92 & 0,006 & \textbf{XXXX} & \textbf{XXXX} \\
литраж   & $-196{,}39$ & 173,67 & $-1{,}13$ & 0,264 & $-546{,}41$ & 153,62 \\
длина    & 121,71 & 54,31 & \textbf{XXXX} & 0,030 & 12,27 & 231,16 \\
const    & $-7961{,}16$ & 11653,58 & $-0{,}68$ & 0,498 & $-31447{,}42$ & 15525,09 \\
\bottomrule
\end{tabular}
\end{center}

\begin{enumerate}[label=\asbuk*), leftmargin=2.2em, itemsep=3pt]
  \item Заполните пропуски в таблице (XXXX).
  \item При каких уровнях значимости регрессия в целом значима?
        Сформулируйте основную и альтернативную гипотезы.
  \item Проверьте значимость коэффициента при переменной «литраж» на уровне
        5\,\%, сформулировав гипотезы.
  \item Исследователь предполагает, что ожидаемая цена иномарки, претерпевшей
        четыре поломки, совпадает с ценой ни разу не ломавшегося отечественного
        автомобиля с теми же остальными характеристиками. Сформулируйте эту
        гипотезу в терминах коэффициентов.
\end{enumerate}

\begin{solution}
$n = 50$, $k = 6$, $n - k = 44$, $t_{0{,}025}(44) = 2{,}015$.

\medskip
\textbf{а)} По \eqref{eq:adj}, \eqref{eq:t}:
\[
  R^2_{\text{adj}} = 1 - \frac{n-1}{n-k}\bigl(1 - R^2\bigr) = 1 - \frac{49}{44}\cdot 0{,}6443 = 0{,}2825,
\]
\[
  t_{\text{длина}} = \frac{121{,}71}{54{,}31} = 2{,}24,
  \qquad
  5102{,}07 \pm 2{,}015\cdot 1749{,}39 = (1576{,}4;\ 8627{,}7).
\]

\medskip
\textbf{б)} $H_0\colon \beta_2 = \beta_3 = \beta_4 = \beta_5 = \beta_6 = 0$,
$H_1$: хотя бы один из них не равен нулю. $P$-значение $F$-статистики равно
0,0013, поэтому регрессия значима на любом уровне $\alpha > 0{,}0013$, в
частности на 1\,\%, 5\,\% и 10\,\%.

\medskip
\textbf{в)} $H_0\colon \beta_5 = 0$, $H_1\colon \beta_5 \ne 0$.
$|t| = 1{,}13 < 2{,}015$ ($P = 0{,}264 > 0{,}05$) — коэффициент незначим.

\medskip
\textbf{г)} Ожидаемая цена иномарки с четырьмя поломками отличается от цены
отечественной машины без поломок (при прочих равных) на
$\beta_4 + 4\beta_3$. Гипотеза: $H_0\colon \beta_4 + 4\beta_3 = 0$.
\end{solution}

% -------------------------------------------------------------------
\problem{6 (задачник ВШЭ, №3.2.7)}

По 30 наблюдениям оценена регрессия (в скобках — стандартные ошибки):
\[
  \widehat{Y}_i = \underset{(2{,}4)}{23{,}5} - \underset{(0{,}2)}{0{,}6}\,X_{2i}
  + \underset{(0{,}7)}{1{,}2}\,X_{3i} + \underset{(0{,}1)}{0{,}9}\,X_{4i}
  - \underset{(1{,}2)}{0{,}2}\,X_{5i},
  \qquad \ESS = 150,\quad R^2 = 0{,}7,
\]
$\ESS$ — сумма квадратов остатков.
\begin{enumerate}[label=\asbuk*), leftmargin=2.2em, itemsep=3pt]
  \item Рассчитайте несмещённую оценку дисперсии случайной составляющей.
  \item Рассчитайте $R^2_{\text{adj}}$.
  \item При исключении каких переменных $R^2_{\text{adj}}$ увеличится?
  \item Постройте 95\,\% доверительный интервал для коэффициента при $X_2$.
  \item На уровне 5\,\% проверьте гипотезу о том, что увеличение $X_4$ на
        единицу при прочих равных увеличивает $Y$ в среднем на 0,5.
  \item Исследователь оценил дополнительную регрессию
        $\widehat{Y}_i = 28{,}4 - 0{,}5(X_{2i} - X_{3i}) + 1{,}1X_{4i}$.
        Какую гипотезу он хотел проверить?
\end{enumerate}

\begin{solution}
$n = 30$, $k = 5$, $n - k = 25$, $t_{0{,}025}(25) = 2{,}060$.

\medskip
\textbf{а)} $s^2 = \ESS/(n-k) = 150/25 = 6$.

\medskip
\textbf{б)} $R^2_{\text{adj}} = 1 - \dfrac{29}{25}\cdot 0{,}3 = 0{,}652$.

\medskip
\textbf{в)} $t$-статистики коэффициентов при $X_2, \dots, X_5$:
$-3$; $1{,}71$; $9$; $-0{,}17$. Исключение одной переменной увеличивает
$R^2_{\text{adj}}$ тогда и только тогда, когда её $|t| < 1$ (для одного
ограничения $F = t^2$, и $R^2_{\text{adj}}$ растёт при $F < 1$). Такая
переменная одна — $X_5$.

\medskip
\textbf{г)} $-0{,}6 \pm 2{,}060\cdot 0{,}2 = (-1{,}01;\ -0{,}19)$.

\medskip
\textbf{д)} $H_0\colon \beta_4 = 0{,}5$, $H_1\colon \beta_4 \ne 0{,}5$:
\[
  t = \frac{0{,}9 - 0{,}5}{0{,}1} = 4 > 2{,}060 ,
\]
гипотеза отвергается.

\medskip
\textbf{е)} В дополнительной регрессии нет $X_5$, а коэффициенты при $X_2$ и
$X_3$ равны по модулю и противоположны по знаку. Значит, проверялась
совместная гипотеза из $q = 2$ ограничений:
$H_0\colon \beta_2 + \beta_3 = 0,\ \beta_5 = 0$. Её проверяют по
\eqref{eq:restr}, сравнив остаточные суммы двух регрессий.
\end{solution}

% -------------------------------------------------------------------
\problem{7 (задачник ВШЭ, №3.2.5)}

По 25 наблюдениям оценено уравнение регрессии количества проданных бутылок
газированной воды $Q_i$ на цену газировки $P_i$ и среднюю температуру дня
$T_i$. В модель включена фиктивная переменная $W_i$, равная единице для
выходных дней и нулю для будних:
\[
  \widehat{Q}_i = \underset{(5{,}7)}{67{,}5} - \underset{(1{,}0)}{2{,}4}\,P_i
  + \underset{(1{,}5)}{2{,}1}\,T_i + \underset{(3{,}2)}{13{,}7}\,W_i,
  \qquad R^2 = 0{,}7 .
\]
\begin{enumerate}[label=\asbuk*), leftmargin=2.2em, itemsep=3pt]
  \item Проверьте значимость регрессии в целом на уровне 1\,\%.
  \item Какие регрессоры значимо влияют на объём продаж (уровень 5\,\%)?
  \item Постройте 95\,\% доверительный интервал для свободного члена.
  \item Пусть вместо $W_i$ в уравнение включена переменная $N_i$, равная нулю
        для выходных и единице для будних дней. Как изменится оценённое
        уравнение?
\end{enumerate}

\begin{solution}
$n = 25$, $k = 4$, $n - k = 21$.

\medskip
\textbf{а)} По \eqref{eq:f}
\[
  F = \frac{R^2}{1-R^2}\cdot\frac{n-k}{k-1} = \frac{0{,}7}{0{,}3}\cdot\frac{21}{3}
    = 16{,}33 > 4{,}87 = F_{0{,}01}(3, 21),
\]
регрессия значима.

\medskip
\textbf{б)} $t$-статистики: $P$: $-2{,}4$; $T$: $1{,}4$; $W$: $4{,}28$.
Критическое значение $t_{0{,}025}(21) = 2{,}080$. Значимы цена и фиктивная
переменная выходного дня; температура незначима.

\medskip
\textbf{в)} $67{,}5 \pm 2{,}080\cdot 5{,}7 = (55{,}6;\ 79{,}4)$.

\medskip
\textbf{г)} $W_i = 1 - N_i$, поэтому
\[
  \widehat{Q}_i = 67{,}5 - 2{,}4P_i + 2{,}1T_i + 13{,}7(1 - N_i)
               = 81{,}2 - 2{,}4P_i + 2{,}1T_i - 13{,}7N_i .
\]
Это та же модель в другой записи: базовой категорией стали выходные. Свободный
член вырос на 13,7, коэффициент при фиктивной переменной сменил знак,
стандартная ошибка при $N_i$ та же (3,2); остальные коэффициенты, их ошибки,
остатки и $R^2$ не изменились. Изменится стандартная ошибка свободного члена
(без ковариаций оценок её не найти).
\end{solution}

% -------------------------------------------------------------------
\problem{8 (задачник ВШЭ, №3.2.10)}

По 40 квартальным наблюдениям оценивается зависимость импорта косметики $M$
от ВВП страны-импортёра $Y$, индекса внутренних цен $\mathit{Pd}$ и индекса
импортных цен $\mathit{Pf}$ в логарифмах. Для сезонности введены фиктивные
переменные $D2$, $D3$, $D4$ (равны единице во 2-м, 3-м и 4-м квартале
соответственно). Получены регрессии ($\ESS$ — сумма квадратов остатков):
\[
\begin{aligned}
  (1)\ \ & \widehat{\ln M} = 7{,}6 + 0{,}8\ln Y + 0{,}1\ln\mathit{Pd} - 0{,}3\ln\mathit{Pf}
           - 0{,}4D2 - 0{,}2D3 - 0{,}2D4, && \ESS = 0{,}37; \\
  (2)\ \ & \widehat{\ln M} = 9{,}3 + 0{,}6\ln Y - 0{,}1(\ln\mathit{Pf} + \ln\mathit{Pd})
           - 0{,}3D2 - 0{,}1D3 - 0{,}3D4, && \ESS = 0{,}48; \\
  (3)\ \ & \widehat{\ln M} = 9{,}0 + 0{,}6\ln Y - 0{,}2(\ln\mathit{Pf} - \ln\mathit{Pd})
           - 0{,}4D2 - 0{,}0D3 - 0{,}3D4, && \ESS = 0{,}40; \\
  (4)\ \ & \widehat{\ln M} = 12{,}9 + 0{,}3\ln\mathit{Pd} + 0{,}1\ln\mathit{Pf}
           - 0{,}2D2 - 0{,}1D3 - 0{,}3D4, && \ESS = 0{,}45; \\
  (5)\ \ & \widehat{\ln M} = 9{,}4 + 0{,}5\ln Y - 0{,}3D2 - 0{,}2D3 - 0{,}3D4, && \ESS = 0{,}42; \\
  (6)\ \ & \widehat{\ln M} = 10{,}1 + 0{,}4\ln Y - 0{,}1\ln\mathit{Pd} + 0{,}1\ln\mathit{Pf}, && \ESS = 1{,}48 .
\end{aligned}
\]
\begin{enumerate}[label=\asbuk*), leftmargin=2.2em, itemsep=3pt]
  \item Проверьте наличие сезонных эффектов на уровне 1\,\%.
  \item На уровне 5\,\% проверьте гипотезу о том, что коэффициент при
        $\ln\mathit{Pd}$ равен по модулю и противоположен по знаку коэффициенту
        при $\ln\mathit{Pf}$, то есть на импорт влияет лишь отношение цен.
\end{enumerate}

\begin{solution}
Модель без ограничений — (1): $k = 7$, $n - k = 33$,
$\ESS_{\mathrm{UR}} = 0{,}37$. Каждая гипотеза — это выбор подходящей модели с
ограничениями.

\medskip
\textbf{а)} Нет сезонности: $H_0\colon \gamma_2 = \gamma_3 = \gamma_4 = 0$
(коэффициенты при $D2$, $D3$, $D4$), $q = 3$. Модель с ограничениями — (6),
в ней нет фиктивных переменных: $\ESS_{\mathrm{R}} = 1{,}48$. По
\eqref{eq:restr}
\[
  F = \frac{\bigl(\ESS_{\mathrm{R}} - \ESS_{\mathrm{UR}}\bigr)/q}
           {\ESS_{\mathrm{UR}}/(n-k)}
    = \frac{\bigl(\ESS_{\mathrm{R}} - \ESS_{\mathrm{UR}}\bigr)/3}
           {\ESS_{\mathrm{UR}}/33},
\]
\[
  F = \frac{(1{,}48 - 0{,}37)/3}{0{,}37/33} = 33{,}0 > 4{,}44 = F_{0{,}01}(3, 33).
\]
Сезонность есть.

\medskip
\textbf{б)} $H_0\colon \beta_{\mathit{Pd}} = -\beta_{\mathit{Pf}}$, $q = 1$.
При этом ограничении
$\beta_{\mathit{Pd}}\ln\mathit{Pd} + \beta_{\mathit{Pf}}\ln\mathit{Pf} =
\beta_{\mathit{Pf}}(\ln\mathit{Pf} - \ln\mathit{Pd})$ — это модель (3),
$\ESS_{\mathrm{R}} = 0{,}40$. (Модель (2) соответствует другому ограничению
$\beta_{\mathit{Pd}} = \beta_{\mathit{Pf}}$.)
\[
  F = \frac{\bigl(\ESS_{\mathrm{R}} - \ESS_{\mathrm{UR}}\bigr)/1}
           {\ESS_{\mathrm{UR}}/33}
    = \frac{0{,}40 - 0{,}37}{0{,}37/33} = 2{,}68 < 4{,}14 = F_{0{,}05}(1, 33).
\]
Гипотеза не отвергается: данные согласуются с тем, что на импорт влияет
отношение цен $\mathit{Pf}/\mathit{Pd}$.
\end{solution}

% -------------------------------------------------------------------
\problem{9 (задачник ВШЭ, №3.2.11)}

По 55 наблюдениям (30 мужчин и 25 женщин) оценена регрессия почасовой
зарплаты $\mathit{EARN}$ на длительность обучения $S$ (годы) и результат ЕГЭ
по математике $\mathit{TEST}$ ($\ESS$ — сумма квадратов остатков):
\begin{center}
\begin{tabular}{l>{$}l<{$}>{$}r<{$}>{$}r<{$}}
\toprule
Выборка & \text{Оценённое уравнение} & \TSS & \ESS \\
\midrule
вся     & \widehat{\mathit{EARN}} = 3744{,}44 - 35{,}54\,S + 3{,}59\,\mathit{TEST}  & 100\,000\,000 & 50\,000\,000 \\
женщины & \widehat{\mathit{EARN}} = 5664{,}05 - 114{,}22\,S - 9{,}05\,\mathit{TEST} &  51\,000\,000 &    550\,000 \\
мужчины & \widehat{\mathit{EARN}} = 2783{,}67 + 112{,}11\,S + 13{,}50\,\mathit{TEST} &  48\,000\,000 &    870\,000 \\
\bottomrule
\end{tabular}
\end{center}
\begin{enumerate}[label=\asbuk*), leftmargin=2.2em, itemsep=3pt]
  \item Найдите коэффициенты уравнения, оценённого по всей выборке:
        \[
          \widehat{\mathit{EARN}} = \widehat{\beta}_1 + \widehat{\beta}_2\,S\cdot\mathit{MALE}
          + \widehat{\beta}_3\,S\cdot\mathit{FEMALE} + \widehat{\beta}_4\,\mathit{FEMALE}
          + \widehat{\beta}_5\,\mathit{TEST}\cdot\mathit{MALE} + \widehat{\beta}_6\,\mathit{TEST},
        \]
        где $\mathit{MALE} = 1$ для мужчин, $\mathit{FEMALE} = 1$ для женщин
        (и 0 иначе).
  \item Есть ли основания считать, что зависимости для мужчин и для женщин
        различаются?
\end{enumerate}

\begin{solution}
\textbf{а)} Подставим значения фиктивных переменных:
\[
\begin{aligned}
  \text{мужчины } (\mathit{MALE} = 1,\ \mathit{FEMALE} = 0):&\quad
    \widehat{\beta}_1 + \widehat{\beta}_2S + (\widehat{\beta}_5 + \widehat{\beta}_6)\,\mathit{TEST}, \\
  \text{женщины } (\mathit{MALE} = 0,\ \mathit{FEMALE} = 1):&\quad
    (\widehat{\beta}_1 + \widehat{\beta}_4) + \widehat{\beta}_3S + \widehat{\beta}_6\,\mathit{TEST}.
\end{aligned}
\]
У каждой группы свои свободный член, наклон по $S$ и наклон по
$\mathit{TEST}$ — по три параметра на группу, всего шесть. Такая модель
подгоняет группы независимо, и её оценки совпадают с оценками отдельных
регрессий. Приравнивая коэффициенты:
\[
\begin{aligned}
  &\widehat{\beta}_1 = 2783{,}67, \qquad \widehat{\beta}_2 = 112{,}11, \qquad
   \widehat{\beta}_5 + \widehat{\beta}_6 = 13{,}50, \\
  &\widehat{\beta}_1 + \widehat{\beta}_4 = 5664{,}05, \qquad \widehat{\beta}_3 = -114{,}22, \qquad
   \widehat{\beta}_6 = -9{,}05,
\end{aligned}
\]
откуда $\widehat{\beta}_4 = 2880{,}38$, $\widehat{\beta}_5 = 22{,}55$.

\medskip
\textbf{б)} Это тест Чоу: $H_0$ — зависимости совпадают ($\beta_2 = \beta_3$,
$\beta_4 = 0$, $\beta_5 = 0$; $q = 3$). Модель с ограничениями — регрессия по
всей выборке, $\ESS_{\mathrm{R}} = 50\,000\,000$. Модель без ограничений — из
п.~а), её остаточная сумма равна сумме остаточных сумм двух отдельных
регрессий: $\ESS_{\mathrm{UR}} = 550\,000 + 870\,000 = 1\,420\,000$, $k = 6$,
$n - k = 49$. По \eqref{eq:restr}
\[
  F = \frac{\bigl(\ESS_{\mathrm{R}} - \ESS_{\mathrm{UR}}\bigr)/q}
           {\ESS_{\mathrm{UR}}/(n-k)}
    = \frac{\bigl(\ESS_{\mathrm{R}} - \ESS_{\mathrm{UR}}\bigr)/3}
           {\ESS_{\mathrm{UR}}/49},
\]
\[
  F = \frac{(50\,000\,000 - 1\,420\,000)/3}{1\,420\,000/49}
    = 558{,}8 > 4{,}21 = F_{0{,}01}(3, 49).
\]
Гипотеза уверенно отвергается: зависимости для мужчин и женщин различаются.
\end{solution}

% -------------------------------------------------------------------
\ifsolutions\clearpage\fi
\problem{10 (задание 2024/25)}

Оценивается модель $Y_t = \alpha + \beta_1X_{1t} + \beta_2X_{2t} +
\beta_3X_{3t} + \varepsilon_t$. Для проверки мультиколлинеарности оценены
вспомогательные регрессии $X_1$ на $X_2, X_3$; $X_2$ на $X_1, X_3$; $X_3$ на
$X_1, X_2$ с коэффициентами детерминации $R_1^2$, $R_2^2$, $R_3^2$. Объясните,
что такое мультиколлинеарность, и сделайте вывод о её наличии для каждого
варианта:
\begin{center}
\begin{tabular}{ccccc}
\toprule
 & Вариант 1 & Вариант 2 & Вариант 3 & Вариант 4 \\
\midrule
$R_1^2$ & 0,80 & 0,10 & 0,73 & 0,90 \\
$R_2^2$ & 0,25 & 0,65 & 0,60 & 0,65 \\
$R_3^2$ & 0,35 & 0,85 & 0,20 & 0,15 \\
\bottomrule
\end{tabular}
\end{center}
Какие есть способы решить проблему мультиколлинеарности?

\begin{solution}
\emph{Мультиколлинеарность} — тесная линейная связь между регрессорами.
\emph{Строгая} (точная линейная зависимость столбцов $\mathbf{X}$) делает
МНК-оценку невозможной: $\mathbf{X}^{\top}\mathbf{X}$ вырождена.
\emph{Нестрогая} оставляет оценки несмещёнными, но увеличивает их дисперсии:
отдельные $t$-статистики малы при высоком $R^2$ и значимом $F$, оценки
неустойчивы, знаки могут быть «неправильными».

Вспомогательный $R_j^2$ показывает, насколько $X_j$ объясняется остальными
регрессорами. По \eqref{eq:vif} $\mathrm{VIF}_j = 1/(1 - R_j^2)$:
\begin{center}
\begin{tabular}{ccccc}
\toprule
 & Вариант 1 & Вариант 2 & Вариант 3 & Вариант 4 \\
\midrule
$\mathrm{VIF}_1$ & 5,00 & 1,11 & 3,70 & \textbf{10,00} \\
$\mathrm{VIF}_2$ & 1,33 & 2,86 & 2,50 & 2,86 \\
$\mathrm{VIF}_3$ & 1,54 & \textbf{6,67} & 1,25 & 1,18 \\
\midrule
Вывод & на границе & есть ($X_3$) & нет & есть, сильная ($X_1$) \\
\bottomrule
\end{tabular}
\end{center}
\begin{itemize}[leftmargin=2.2em, itemsep=2pt]
  \item \emph{Вариант 1}: $\mathrm{VIF}_1 = 5$ — ровно на пороге 5;
        умеренная мультиколлинеарность, связанная с $X_1$; по порогу 10 —
        нет.
  \item \emph{Вариант 2}: $\mathrm{VIF}_3 = 6{,}67 > 5$ — мультиколлинеарность
        есть, $X_3$ сильно связан с остальными (но $< 10$, не критическая).
  \item \emph{Вариант 3}: все VIF $< 5$ — мультиколлинеарности нет.
  \item \emph{Вариант 4}: $\mathrm{VIF}_1 = 10$ — серьёзная
        мультиколлинеарность, дисперсия $\widehat{\beta}_1$ в 10 раз больше,
        чем при несвязанных регрессорах.
\end{itemize}

\textbf{Способы решения.}
\begin{enumerate}[label=\arabic*), leftmargin=2.2em, itemsep=2pt]
  \item Исключить одну из коррелирующих переменных (ту, что хуже объясняет
        $Y$ в парной регрессии). Цена — смещение оставшихся коэффициентов,
        если исключённая переменная на самом деле влияет на $Y$.
  \item Получить дополнительные данные или новую выборку (например, перейти
        от годовых данных к квартальным).
  \item Использовать априорную информацию о параметрах: наложить ограничение
        (например, $\beta_1 + \beta_2 = 1$) или объединить коррелирующие
        переменные в одну.
  \item Преобразовать данные: перейти к приростам или к относительным
        величинам (разделить на одну из коррелирующих переменных).
\end{enumerate}
\end{solution}

% -------------------------------------------------------------------
\problem{11 (задачник ВШЭ, №4.1.3)}

По наблюдениям за $Y$, $X_2$ и $X_3$ оценена корреляционная матрица:
\begin{center}
\begin{tabular}{c|ccc}
      & $X_2$  & $X_3$  & $Y$ \\ \hline
$X_2$ & 1,0    & $-0{,}9$ & $-0{,}6$ \\
$X_3$ & $-0{,}9$ & 1,0    & 0,7 \\
$Y$   & $-0{,}6$ & 0,7    & 1,0
\end{tabular}
\end{center}
Найдите $\mathrm{VIF}(\widehat{\beta}_3)$ для регрессии
$Y_i = \beta_1 + \beta_2X_{2i} + \beta_3X_{3i} + \varepsilon_i$.

\begin{solution}
Регрессоров два, поэтому вспомогательная регрессия $X_3$ на константу и $X_2$
— парная, и её $R^2$ равен квадрату коэффициента корреляции:
$R_3^2 = r^2(X_2, X_3) = 0{,}81$. Корреляции с $Y$ для VIF не нужны.
Мультиколлинеарность заметная (VIF $> 5$), но ниже порога 10:
\[
  \mathrm{VIF}(\widehat{\beta}_3) = \frac{1}{1 - 0{,}81} = 5{,}26 .
\]
\end{solution}

% -------------------------------------------------------------------
\ifsolutions\clearpage\fi
\problem{12 (задачник ВШЭ, №4.1.11)}

Для уравнения $Y_i = \beta_1 + \beta_2X_{2i} + \beta_3X_{3i} + \beta_4X_{4i} +
\varepsilon_i$ получены корреляционная матрица регрессоров и значения VIF:
\begin{center}
\begin{tabular}{c|ccc}
      & $X_2$ & $X_3$ & $X_4$ \\ \hline
$X_2$ & 1 & & \\
$X_3$ & $-0{,}4934$ & 1 & \\
$X_4$ & $-0{,}6144$ & 0,9018 & 1
\end{tabular}
\qquad
\begin{tabular}{cc}
\toprule
Переменная & VIF \\
\midrule
$X_4$ & 2,53 \\
$X_3$ & 2,08 \\
$X_2$ & 1,62 \\
\bottomrule
\end{tabular}
\end{center}
Объясните, почему эти VIF и корреляции не могут относиться к одним и тем же
данным.

\begin{solution}
$R_4^2$ — коэффициент детерминации регрессии $X_4$ на $X_2$ и $X_3$. Добавление
регрессора не уменьшает $R^2$, поэтому $R_4^2$ не меньше коэффициента
детерминации парной регрессии $X_4$ на одно $X_3$, а он равен
$r^2(X_3, X_4)$:
\[
  R_4^2 \ \ge\ r^2(X_3, X_4) = 0{,}9018^2 = 0{,}813
  \quad\Longrightarrow\quad
  \mathrm{VIF}_4 = \frac{1}{1 - R_4^2} \ \ge\ \frac{1}{1 - 0{,}813} = 5{,}35 .
\]
Так же $\mathrm{VIF}_3 \ge 5{,}35$. А в таблице $\mathrm{VIF}_4 = 2{,}53$ и
$\mathrm{VIF}_3 = 2{,}08$ — меньше нижней границы. Значит, VIF и корреляции
относятся к разным данным (или в расчётах ошибка).
\end{solution}
